\section{Scope, Method, and Historical Coordinates}

\subsection{Identity and corpus boundary}

Subject: the Latin biblical tradition conventionally called the Vulgate,
treated as the layered object the body of this account distinguishes: the
pre-Jerome Old Latin (\work{Vetus Latina}) background; Jerome's Gospel
revision (382--384); the anonymous revision of the remaining New
Testament (Rome, c.~393--410); Jerome's three Psalters and his
Hebrew-based Old Testament (c.~390--405/6) with the Aramaic-based Tobit
and Judith; the untouched Old Latin deuterocanonical books; the medieval
transmission and its editions (Cassiodorus, Amiatinus, Alcuin, Theodulf,
Paris Bibles, \textit{correctoria}); the printed and official editions
(Gutenberg; Sixtine 1590; Sixto-Clementine 1592); the modern critical
editions (Oxford, Rome, Stuttgart); and the \work{Nova Vulgata} (1979;
\textit{editio typica altera} 1986). ``Vulgate'' is used of the received
collection anachronistically but conventionally; where a specific edition
matters, it is named. ``Old Latin''/\work{Vetus Latina} is a modern
collective label, not one ancient edition. \textit{Editio vulgata} in
fourth-century usage denoted the common Septuagint-based text.
``Authentic'' bears only the juridical sense defined at Trent as glossed
by \work{Divino afflante Spiritu}. ``Typical edition'' is the normative
edition for a defined ecclesial use, not a claim of antiquity.
``Hebrew verity'' (\textit{Hebraica veritas}) names Jerome's program and
polemic, not a claim that his exemplars equal the medieval Masoretic
Text.

\subsection{Coverage and exclusions}

Chronological center: the second through seventeenth centuries, with a
bounded modern extension (critical editions, \work{Nova Vulgata},
liturgical instructions, Catechism) checked through 2026-07-24.
Geographic center: Rome, Bethlehem, North Africa, and the Latin West.
This account creates no translation of any ancient, biblical, or official
text; quotations follow the identified public-domain English versions in
the References, and copyrighted works (notably Houghton's guide, cited
CC~BY-NC-ND) are paraphrased or cited, never reproduced at length.
Excluded: verse-level textual commentary; any manuscript census or
stemma; complete histories of Latin exegesis, vernacular Bibles, or
liturgical translation; any adjudication of canon or inspiration (acts
are reported, not judged); any election of one edition as ``the real
Vulgate''; and any claim that presently binding liturgical law is stated
here beyond the cited instructions' own dates.

\subsection{Evidence classes}

\begin{evidencekey}
\evidenceclass{D}{Documents and artifacts: manuscripts and printed
artifacts cited at catalog level (Codex Amiatinus, Tours and Theodulf
Bibles, the Gutenberg Bible, the Hetzenauer 1914 Clementine printing,
page image inspected).}
\evidenceclass{A}{Contemporary participant documents read at their own
date and purpose: Jerome's prefaces and letters, Augustine's letters,
Rufinus's apology, \work{Primum quaeritur}.}
\evidenceclass{T}{Later tradition and reception: Bede on Ceolfrith,
Cassiodorus's descriptions, Bellarmine's preface as reported, the
Douay--Rheims enterprise.}
\evidenceclass{M}{Official ecclesial acts, each limited to its object:
Trent Session IV; \work{Divino afflante Spiritu} 21; \work{Dei Verbum}
22; \work{Scripturarum thesaurus}; \work{Nova Vulgata} praefatio;
\work{Liturgiam authenticam}; CCC 131.}
\evidenceclass{K}{Named modern critical reconstruction and survey
literature: Houghton 2016 above all; De Bruyne; the edition projects;
dated Catholic Encyclopedia articles used as identified period
witnesses.}
\evidenceclass{S}{Source-grounded project synthesis, never attributed to
a witness.}
\end{evidencekey}

\subsection{Method, uncertainties, and negative results}

Participant prefaces establish what their authors claimed, at their date,
in their polemical situation; no preface is treated as a neutral
administrative record. Official acts establish the status they confer and
nothing about earlier history. Modern reconstructions are attributed by
name. Material uncertainties preserved in place: the exact content of
Damasus's request; the identity of the rest-of-New-Testament reviser
(Pelagius and ``Rufinus the Syrian'' remain hypotheses); the identity of
the extant Roman Psalter with Jerome's Rome revision (doubted since De
Bruyne); the per-book dates of the Hebrew-based translations (only a
sequence is recoverable); the sources and procedure behind Tobit and
Judith; what Jerome meant by his Nicene-Council remark about Judith (no
surviving Nicene canon lists Scripture); the textual identities of
Cassiodorus's three Bibles; the promulgation status of \textit{Aeternus
ille} and the full motives for the Sixtine withdrawal, where published
difference-counts (c.~2,000--5,000) are reported as a spread, none
endorsed. Bounded negative results: no commission instrument from
Damasus survives; no ancient source names the reviser of Acts, the
Epistles, and Revelation; the Catechism's article on Scripture (CCC
101--141) was checked and does not treat the Vulgate by name. These are
correctable search results, not proofs of absence. Mutable web sources
were used as checked on 2026-07-24; the quoted witnesses were re-verified
on 2026-07-25.

\subsection{Rights and review state}

Project prose is original and intended for CC~BY~4.0; ancient and
official texts quoted follow public-domain translations identified in
the References (Edgecomb's preface translations, released to the public
domain; NPNF; Dods; Waterworth's Trent); Vatican web texts are quoted
within ordinary scholarly limits with citation; Houghton 2016
(CC~BY-NC-ND) is summarized without adaptation. This publication has
received internal editorial checking only. No independent patristic,
Latin-biblical, medieval, early-modern, theological, canonical, or
rights review is claimed, and no ecclesiastical review or approval is
claimed or implied.

\subsection{Orientation timeline}

Dates before print are approximate unless a witness supplies them;
``convention'' marks customary datings this account does not assert as
established.

\begin{accounttimeline}
\timelineevent{2nd--3rd c.}{Latin translations of Scripture from Greek in
circulation}{Quotations in African writers; later manuscript tradition
(D/K)}{Secure as fact; origins reconstructed}
\timelineevent{382--384}{Jerome revises the four Gospels at Rome under
Damasus}{Preface \textit{Novum opus} (A); \work{Letter} 27 (A)}{Secure;
commission's exact terms lost}
\timelineevent{384}{Damasus dies; Jerome leaves Rome 385, Bethlehem
386}{Jerome's letters (A); chronology conventional}{Secure in outline}
\timelineevent{c.~386--390}{Hexaplaric revisions of Old Testament books;
Gallican Psalter}{Psalter preface (A); \work{Letter} 112 (A)}{Secure;
most revisions lost}
\timelineevent{c.~390--405/6}{Old Testament translated from the Hebrew;
Tobit and Judith from Aramaic}{Prefaces (A); sequence only (K)}{Secure as
program; per-book dates open}
\timelineevent{c.~393--410}{Anonymous revision of Acts, Epistles,
Revelation at Rome}{\work{Primum quaeritur} (A); dating bounds
(K)}{Secure; reviser unknown}
\timelineevent{403--405}{Augustine--Jerome exchange; Oea incident;
Rufinus's attack}{\work{Letters} 71, 112, 82 (A); \work{Apology} 2
(A)}{Secure witnesses}
\timelineevent{419/420}{Jerome dies at Bethlehem}{Chronicle tradition
(T)}{Year conventional}
\timelineevent{c.~540s--580s}{Cassiodorus's Bibles at Vivarium}{
\work{Institutiones} as reported (T/K)}{Secure; textual identities
disputed}
\timelineevent{before 716}{Codex Amiatinus written at
Wearmouth--Jarrow}{Bede (T); the manuscript (D)}{Secure}
\timelineevent{796--800}{Alcuin's revision at Tours; Theodulf's editions
at Orl\'eans}{Alcuin's letters; surviving Bibles (D/K)}{Secure; no
official imperial text}
\timelineevent{c.~1230}{Paris Bibles standardize order and modern
chapters}{Surviving Bibles (D/K)}{Secure; Langton attribution
qualified}
\timelineevent{1454--1456}{Gutenberg's 42-line Bible printed at
Mainz}{Surviving copies; Ransom Center description (D)}{Secure}
\timelineevent{1546}{Trent: canon decree; ``old and vulgate edition''
authentic; corrected printing ordered}{Session IV decrees (M)}{Secure}
\timelineevent{1582/1609--10}{Rheims New Testament; Douay Old Testament
from the ``authentical Latin''}{Title pages; survey literature
(T/K)}{Secure}
\timelineevent{1590}{Sixtine Vulgate printed; withdrawn after Sixtus V's
death}{Editions; Bellarmine's account as reported (D/T)}{Secure sequence;
motives disputed}
\timelineevent{1592}{Sixto-Clementine Vulgate; standard text to
1979}{The edition (D); reception (K)}{Secure}
\timelineevent{1889--1954}{Oxford Vulgate New Testament
(Wordsworth--White)}{Edition record (K)}{Secure}
\timelineevent{1907; 1926--1995}{Benedictine commission; Roman Vulgate
Old Testament}{Gasquet 1912 (K); edition record (K)}{Secure}
\timelineevent{1943}{\work{Divino afflante Spiritu} glosses Trent's
authenticity as juridical}{Encyclical text (M)}{Secure}
\timelineevent{1945}{Vetus Latina Institute founded at Beuron}{Project
record (K)}{Secure}
\timelineevent{1965}{\work{Dei Verbum} 22; Nova Vulgata commission
established}{Conciliar text (M); \work{Praefatio} (M)}{Secure}
\timelineevent{1969/2007}{Stuttgart Vulgate; fifth edition
(Weber--Gryson)}{Edition record (K)}{Secure}
\timelineevent{1979/1986}{\work{Scripturarum thesaurus} promulgates the
Nova Vulgata; second typical edition}{Constitution and \work{Praefatio}
(M)}{Secure}
\timelineevent{2001}{\work{Liturgiam authenticam} assigns the Nova
Vulgata its reference role in liturgical translation}{Instruction text
(M)}{Secure as act of its date}
\end{accounttimeline}
